% !TEX root = ../../../main.tex
% Block card: detection.signal_threshold -- "Signal threshold to spans" (Part I, the blocks chapter).
\cardfield{What it does}
Marks every stretch where the signal is above, or below, a level you type.

\cardfield{When to use it}
When you know the size of the events you want: ``every stretch where the trace falls more
than 5 millivolts under its baseline''. It is the simplest detector in the app. Put
\blockref{preprocessing.detrend} in front so the level is measured from the baseline; the
ready-made template \texttt{signal\_threshold} does that and looks for drops below $-5$\,mV.

\cardfield{What goes in, what comes out}
In goes a stretch of one channel. Out come the stretches past the level, each scored by how
far past it they reach.

\cardfield{Settings that matter}
\emph{threshold} is the level; you can drag it on the trace. \emph{direction}: \texttt{above}
(the default) keeps stretches over the level, \texttt{below} keeps stretches under it (give a
negative level for drops), \texttt{either} keeps stretches further from zero than the level on
either side. \emph{min\_duration\_s} and \emph{merge\_gap\_s} tidy the result: drop very short
stretches, join ones that are nearly touching.

\cardfield{Watch out for}
The level is in the unit the recording is \emph{stored} in, not always millivolts. The M2
recordings are stored in volts, so $-5$\,mV is typed as $-0.005$; on a recording stored in
millivolts the same $-0.005$ is a thousand-fold smaller cut and catches noise. The default level,
0 with \texttt{above}, keeps every stretch above zero: always set it.
